\section{总结与延伸}

\subsection{核心要点回顾}

本讲详细推导了 DDPM 的训练和采样算法，核心内容可以总结为以下几点：

\begin{enumerate}
    \item \textbf{前向过程}是预定义的高斯噪声注入过程，具有闭式的跳跃分布 $q(x_t|x_0)$，允许一步采样
    \item \textbf{ELBO 分解}将训练目标转化为在每个时间步匹配真实后验转移分布
    \item \textbf{真实后验} $q(x_{t-1}|x_t, x_0)$ 是高斯分布，方差可直接计算，均值依赖 $x_0$
    \item \textbf{三种参数化}：均值预测、$x_0$ 预测、噪声预测在数学上等价，实践中\textbf{噪声预测}最稳定
    \item \textbf{简化训练目标}：$L_{\text{simple}} = \mathbb{E}\|\epsilon - \epsilon_\theta(x_t, t)\|^2$，优雅而高效
    \item \textbf{表达能力}来自反向过程中神经网络引入的非线性：每步高斯但整体非高斯
    \item \textbf{噪声预测}与 score-based 模型有深刻联系，为统一理论框架铺平道路
\end{enumerate}

\subsection{后续课程预告}

讲者提到后续将讨论：
\begin{itemize}
    \item 噪声预测与 score matching 的正式联系
    \item Score SDE 框架：将离散时间扩散推广到连续时间
    \item 不同的前向过程定义方式
    \item 更快的采样方法（DDIM 等）
\end{itemize}

\subsection{拓展阅读}

\begin{itemize}
    \item \textbf{DDPM 原论文}：Ho, Jain, Abbeel. ``Denoising Diffusion Probabilistic Models.'' NeurIPS 2020.
    \item \textbf{Improved DDPM}：Nichol, Dhariwal. ``Improved Denoising Diffusion Probabilistic Models.'' ICML 2021. 探索了学习方差的方案。
    \item \textbf{Score SDE}：Song et al. ``Score-Based Generative Modeling through Stochastic Differential Equations.'' ICLR 2021. 统一了 DDPM 和 score-based 模型。
    \item \textbf{DDIM}：Song, Meng, Ermon. ``Denoising Diffusion Implicit Models.'' ICLR 2021. 提出了确定性采样方法，大幅减少采样步数。
    \item \textbf{P2 Weighting}：Choi et al. ``Perception Prioritized Training of Diffusion Models.'' CVPR 2022. 探索了时间步权重的自适应策略。
    \item \textbf{v-prediction}：Salimans, Ho. ``Progressive Distillation for Fast Sampling of Diffusion Models.'' ICLR 2022. 提出了结合 $x_0$ 预测和 $\epsilon$ 预测优点的 v-prediction 参数化。
\end{itemize}

\end{document}
